\chapter{Méthodologie et Architecture}

\section*{Introduction}

L'évolution rapide de l'intelligence artificielle dans les systèmes de sécurité a engendré des défis sans précédent en matière de robustesse et de fiabilité. Dans ce contexte, ce mémoire s'est engagé à explorer les différentes stratégies de sécurisation des modèles d'IA face aux attaques adversariales dans le domaine du maintien de l'ordre. Le chapitre précédent a déjà permis d'analyser les vulnérabilités des systèmes d'IA, d'examiner les menaces existantes et d'évaluer les risques associés aux attaques adversariales. Dans ce chapitre, nous décrirons en détail les étapes de l'implémentation de notre architecture de sécurisation, en mettant en lumière les outils, les paramètres et les méthodes utilisés.

\section{Démarche méthodologique}

La démarche méthodologique de cette étude repose sur une approche rigoureuse visant à explorer et à mettre en œuvre des stratégies de sécurisation des modèles d'IA dans les systèmes de maintien de l'ordre. Cette section détaille les étapes suivies pour atteindre les objectifs fixés, fournissant ainsi un aperçu du processus de recherche.

\subsection{Collecte des Données}

La collecte des données repose sur une approche en deux phases complémentaires. Dans un premier temps, nous avons procédé à l'identification des méthodes et outils pertinents pouvant intervenir dans la sécurisation des modèles d'IA dans le domaine du maintien de l'ordre. Cette phase exploratoire a nécessité une revue de la littérature scientifique sur les techniques de défense adversariale \cite{madry2017towards}, l'analyse des frameworks de sécurité existants, ainsi qu'une étude comparative des bibliothèques disponibles pour l'implémentation de mécanismes de protection.

Dans un second temps, nous nous sommes concentrés sur l'identification et l'étude d'un cas pratique nous permettant d'illustrer la méthode de sécurisation proposée. Nous avons sélectionné un problème de classification d'images représentatif des applications de maintien de l'ordre, tel que la détection d'objets suspects. Les données ont été recueillies à partir de sources fiables et publiques, garantissant une représentation réaliste des défis rencontrés dans le déploiement opérationnel de tels systèmes.

\subsection{Modélisation d'une Architecture de Sécurité}

Nous avons procédé à la modélisation d'une architecture de sécurité modulaire s'appuyant sur les principes de défense en profondeur adaptés aux spécificités des systèmes d'apprentissage automatique. L'architecture proposée s'articule autour de plusieurs couches de protection complémentaires couvrant l'ensemble du cycle de vie des modèles d'IA : sécurisation des données d'entraînement, renforcement du modèle via adversarial training, et protection au moment de l'inférence avec détection d'exemples adversariaux. Des composants de monitoring et d'audit permettent une surveillance continue du comportement du modèle en production.

\subsection{Mise en place d'un cas expérimental}

La mise en place du cas expérimental a nécessité la conception d'un environnement de test complet et reproductible. Nous avons développé un pipeline expérimental comprenant : un modèle de référence (baseline), un module de génération d'attaques adversariales (FGSM, PGD, C\&W), et des mécanismes d'évaluation automatisés mesurant l'accuracy, la robustesse, et le temps d'inférence. Le modèle a été déployé dans un environnement conteneurisé avec Docker, exposant une API REST pour simuler des conditions opérationnelles réalistes.

\subsection{Simulation et Évaluation des Résultats}

Les scénarios d'attaque et de défense ont été soumis à des simulations extensives selon un protocole d'évaluation rigoureux. L'évaluation s'articule autour de trois axes principaux : la robustesse adversariale face à différents budgets de perturbation, l'analyse du compromis entre robustesse et accuracy \cite{tsipras2018robustness}, et les performances opérationnelles (latence, ressources, montée en charge).

Une analyse comparative a été conduite pour démontrer la plus-value de notre méthode de sécurisation face aux mécanismes classiques et aux architectures existantes. Nous avons comparé notre approche avec des modèles non sécurisés (baseline), des modèles utilisant uniquement des techniques de défense standard, et des architectures de sécurisation proposées dans la littérature récente. Cette comparaison permet de quantifier les gains en robustesse, d'évaluer le surcoût en performance, et de positionner notre contribution par rapport à l'état de l'art. Cette démarche méthodologique vise à garantir la validité et la fiabilité des résultats obtenus, offrant ainsi une base solide pour l'analyse des stratégies de sécurisation des modèles d'IA au sein des systèmes de maintien de l'ordre.

\section{Méthode de sécurisation des modèles IA dans les systèmes de sécurité}

Face aux menaces qui ont été présentées dans le chapitre précédent, nous proposons une architecture de sécurisation globale et modulaire pour les modèles d'IA déployés dans les systèmes de maintien de l'ordre. Cette architecture s'appuie sur les concepts de défense en profondeur de la cybersécurité classique, mais les adapte aux particularités des systèmes d'apprentissage automatique. Notre approche englobe tout le cycle de vie des modèles IA, de la collecte des données à la mise en production en passant par l'entraînement, la validation et le monitoring continu.

\section{Outils et Technologies}

L'implémentation de notre architecture de sécurisation nécessite un ensemble d'outils et de technologies soigneusement sélectionnés pour leur maturité, leur performance et leur compatibilité avec les contraintes opérationnelles des systèmes de maintien de l'ordre. Cette section présente la stack technologique retenue et justifie nos choix à la lumière des exigences de sécurité, de performance et de reproductibilité.

\subsection{Environnement de développement}

\paragraph{Python 3.8+}
Python constitue le langage principal pour l'implémentation des modèles et des pipelines de sécurité. Ce choix s'appuie sur l'écosystème riche en bibliothèques d'apprentissage automatique et de sécurité dont dispose Python. La maturité de l'écosystème Python pour le machine learning, avec des bibliothèques comme NumPy, scikit-learn et PyTorch, facilite l'intégration de techniques de sécurisation avancées tout en maintenant une base de code maintenable.

\paragraph{PyTorch 2.0+}
PyTorch a été privilégié comme framework de deep learning pour sa flexibilité et son approche dynamique de construction de graphes computationnels \cite{paszke2019pytorch}. Cette caractéristique permet une implémentation facile des techniques d'adversarial training et un contrôle fin du processus d'entraînement, essentiel pour l'intégration de mécanismes de défense personnalisés. La communauté active de PyTorch et la disponibilité d'implémentations de référence pour les techniques de robustesse adversariale ont également motivé ce choix.

\paragraph{Google Colab / Jupyter Notebook}
L'environnement de développement interactif Jupyter facilite l'expérimentation itérative et la documentation du code. Google Colab offre en complément un accès à des ressources GPU gratuites, permettant l'entraînement de modèles sans infrastructure locale coûteuse. Cette approche favorise la reproductibilité des expériences et facilite le partage des résultats avec la communauté scientifique.

\subsection{Bibliothèques de sécurité et de robustesse}

\paragraph{Adversarial Robustness Toolbox (ART)}
L'Adversarial Robustness Toolbox, développée par IBM Research, constitue une bibliothèque de référence dans le domaine de la robustesse adversariale \cite{nicolae2018adversarial}. ART fournit des implémentations standardisées et bien documentées pour :
\begin{itemize}
    \item La génération d'attaques adversariales (FGSM, PGD, C\&W, DeepFool, etc.)
    \item Les méthodes de défense (adversarial training, feature squeezing, defensive distillation)
    \item Les métriques de robustesse permettant d'évaluer quantitativement la résistance des modèles
\end{itemize}
L'utilisation d'ART garantit la conformité avec les implémentations de référence utilisées dans la littérature scientifique, facilitant ainsi la comparaison de nos résultats avec l'état de l'art.

\paragraph{CleverHans}
CleverHans, développée initialement par des chercheurs de Google Brain, est une bibliothèque complémentaire pour les attaques adversariales. Elle est particulièrement réputée pour ses implémentations efficaces des méthodes FGSM et PGD, et pour ses tutoriaux détaillés facilitant la compréhension des mécanismes d'attaque. Nous l'utilisons en complément d'ART pour valider la robustesse de nos défenses face à des implémentations alternatives.

\paragraph{Foolbox}
Foolbox est un framework Python agnostique aux bibliothèques de deep learning, permettant d'évaluer la robustesse des modèles face à diverses attaques adversariales. Son support natif de PyTorch et son interface intuitive en font un outil précieux pour l'évaluation comparative de la robustesse. Foolbox se distingue par sa capacité à générer des perturbations adversariales minimales, fournissant ainsi des insights sur la distance réelle aux frontières de décision.

\subsection{Outils de gestion des données}

\paragraph{DVC (Data Version Control)}
DVC est un système de versioning spécialisé pour les datasets et les modèles d'apprentissage automatique. Il permet de :
\begin{itemize}
    \item Tracer l'évolution des données d'entraînement et des artefacts de modèles
    \item Assurer la reproductibilité complète des expériences
    \item Gérer efficacement les datasets volumineux sans saturer les dépôts Git
    \item Créer des pipelines reproductibles reliant données, code et modèles
\end{itemize}
L'utilisation de DVC est essentielle pour garantir la traçabilité et l'auditabilité des expériences, des exigences critiques dans le contexte des systèmes de sécurité.

\paragraph{Great Expectations}
Great Expectations est un framework de validation de données permettant de définir des attentes explicites sur les propriétés statistiques des datasets. Il permet de détecter les anomalies dans les données d'entraînement, les dérives de distribution, et les tentatives potentielles d'empoisonnement. L'intégration de Great Expectations dans notre pipeline de données fournit une première ligne de défense contre les attaques visant l'intégrité des données.

\subsection{Infrastructure de déploiement}

\paragraph{Docker}
La conteneurisation avec Docker garantit un déploiement isolé, reproductible et sécurisé des modèles. Docker permet d'encapsuler l'ensemble des dépendances et de la configuration, réduisant ainsi les risques liés aux différences d'environnement entre développement et production. L'isolation fournie par les conteneurs constitue également une mesure de sécurité limitant l'impact potentiel d'une compromission.

\paragraph{FastAPI}
FastAPI est un framework web moderne pour la création d'APIs REST performantes en Python. Sa validation automatique des données d'entrée basée sur les type hints Python et Pydantic fournit une protection contre les entrées malformées. FastAPI supporte nativement les opérations asynchrones, permettant d'atteindre des performances élevées tout en maintenant une surface d'attaque réduite grâce à sa conception sécurisée par défaut.

\paragraph{Nginx}
Nginx agit comme reverse proxy, fournissant plusieurs fonctionnalités de sécurité essentielles :
\begin{itemize}
    \item Rate limiting pour prévenir les attaques par déni de service
    \item Load balancing pour distribuer la charge entre plusieurs instances du modèle
    \item Terminaison TLS/SSL pour chiffrer les communications
    \item Filtrage des requêtes malformées avant qu'elles n'atteignent l'application
\end{itemize}

\subsection{Monitoring et observabilité}

\paragraph{Prometheus + Grafana}
Prometheus et Grafana constituent une stack de monitoring permettant de collecter et visualiser les métriques de performance et de sécurité en temps réel. Nous surveillons notamment :
\begin{itemize}
    \item Les métriques de performance du modèle (latence, throughput, accuracy)
    \item Les indicateurs de sécurité (taux de détection d'anomalies, distribution des confidences)
    \item Les métriques d'infrastructure (utilisation CPU/GPU, mémoire, réseau)
\end{itemize}
Cette observabilité continue permet de détecter rapidement les comportements anormaux pouvant indiquer une attaque en cours.

\paragraph{ELK Stack (Elasticsearch, Logstash, Kibana)}
La stack ELK fournit une solution de logging centralisé assurant la traçabilité complète des prédictions et des événements de sécurité. Chaque requête de prédiction est journalisée avec ses métadonnées, permettant des analyses post-mortem en cas d'incident et facilitant les audits de conformité. L'indexation performante d'Elasticsearch permet des recherches rapides dans les logs pour identifier des patterns d'attaque.

\subsection{Justification des choix technologiques}

Le choix de cette stack technologique repose sur plusieurs critères complémentaires : la maturité et la fiabilité des outils, leur compatibilité et interopérabilité via des standards ouverts, leur capacité à atteindre les objectifs de performance (< 200ms de latence), l'accessibilité grâce à des solutions open-source, la scalabilité de l'architecture conteneurisée, et enfin leur adoption dans la recherche scientifique garantissant la comparabilité avec l'état de l'art. Cette stack constitue ainsi une fondation solide alliant rigueur scientifique, performance opérationnelle et pragmatisme économique.

\section*{Conclusion}

Ce chapitre a présenté la méthodologie et l'architecture technologique sous-tendant notre approche de sécurisation des modèles d'IA pour les systèmes de maintien de l'ordre. La démarche méthodologique rigoureuse, articulée autour de la collecte de données, de la modélisation, de l'expérimentation et de l'évaluation, garantit la validité scientifique de nos travaux. La stack technologique sélectionnée, combinant des outils matures de l'écosystème Python avec des bibliothèques spécialisées en robustesse adversariale, fournit les fondations nécessaires à l'implémentation concrète de notre architecture de sécurisation.

Le chapitre suivant présentera l'architecture détaillée de notre solution de sécurisation, décrivant les mécanismes de défense implémentés à chaque couche et leur intégration dans un système cohérent et opérationnel.